
\medskip
\phantomsection\label{par:sga3-VIB-10.14}%
\noindent\textbf{10.14.---}
Statements \hyperref[lem:sga3-VIB-10.2]{10.2} and
\hyperref[lem:sga3-VIB-10.3]{10.3}, concerning the category
of \(S\)-groups of finite presentation, extend to the category
of pairs consisting of an \(S\)-group of finite presentation
and an \(S\)-scheme of finite presentation with operator group
\(G\). More precisely, in the situation recalled at the
beginning of \hyperref[par:sga3-VIB-10.0]{10.0}:

\begin{enumerate}
\item[\textup{(i)}]\phantomsection\label{concl:sga3-VIB-10.14-i}
  let \(j\in I\), and let \(G_j\) and \(G'_j\) be two
  \(S_j\)-groups of finite presentation; let \(H_j\)
  \textup{(respectively, \(H'_j\))} be an \(S_j\)-scheme of
  finite presentation with operator group \(G_j\)
  \textup{(respectively, \(G'_j\))}. For \(i\in I\),
  \(i\geq j\), put
  \[
    G_i=G_j\times_{S_j}S_i,
    \qquad
    G=G_j\times_{S_j}S,
  \]
  and define \(G'_i\), \(G'\), \(H_i\), \(H\), \(H'_i\),
  and \(H'\) similarly. Denote by
  \[
    \operatorname{Dihom}_{S\text{-}\mathrm{gr.}}
      \bigl((G,H),(G',H')\bigr)
  \]
  the set of di-morphisms of \(S\)-groups and \(S\)-schemes
  with operator groups from the pair \((G,H)\) to the pair
  \((G',H')\). Then the canonical map
  \phantomsection\label{form:sga3-VIB-10.14-Dihom-limit}%
  \[
    \varinjlim_{i\geq j}
      \operatorname{Dihom}_{S_i\text{-}\mathrm{gr.}}
        \bigl((G_i,H_i),(G'_i,H'_i)\bigr)
      \longrightarrow
      \operatorname{Dihom}_{S\text{-}\mathrm{gr.}}
        \bigl((G,H),(G',H')\bigr)
  \]
  is bijective.

\item[\textup{(ii)}]\phantomsection\label{concl:sga3-VIB-10.14-ii}
  let \(G\) be an \(S\)-group of finite presentation and
  \(H\) an \(S\)-scheme of finite presentation with operator
  group \(G\). Then there are an index \(j\in I\), an
  \(S_j\)-group \(G_j\) of finite presentation, an
  \(S_j\)-scheme \(H_j\) of finite presentation with operator
  group \(G_j\), and a di-isomorphism of \(S\)-groups and
  \(S\)-schemes with operator groups from
  \[
    \bigl(G_j\times_{S_j}S,\ H_j\times_{S_j}S\bigr)
  \]
  onto \((G,H)\).%
  \begingroup
  \renewcommand{\thefootnote}{\fnsymbol{footnote}}%
  \footnote[1]{\label{note:sga3-VIB-10.14-GH}
  Translator's note: the French re-edition prints \((G,X)\)
  here. Since this assertion begins with the pair \((G,H)\)
  and \(X\) is not defined in \SGARefLink{par:sga3-VIB-10.14}{10.14}, \(H\) is restored.}%
  \endgroup
\end{enumerate}

\medskip
\phantomsection\label{def:sga3-VIB-10.15}%
\noindent\textbf{Definition 10.15.---}%
\footnote{\label{note:sga3-VIB-102}Editor's note: the original
has been corrected by distinguishing the notions of a
``formally homogeneous'' object and a ``homogeneous'' object;
see IV, \S\S \SGARefLink{sga3:iv:subsection:5:1}{5.1} and \SGARefLink{sga3:iv:subsection:6:7}{6.7}.}
\emph{Let \(\mathcal T\) be a topology on
\((\mathbf{Sch}/S)\) coarser than the canonical topology.
Given an \(S\)-group scheme \(G\) and an \(S\)-scheme \(X\)
with operator group \(G\), one says that \(X\) is a
\emph{formally homogeneous space under \(G\)}
\textup{(relative to the topology \(\mathcal T\))} if the
morphism}
\phantomsection\label{form:sga3-VIB-10.15-Phi}%
\[
  \Phi:G\times_S X\longrightarrow X\times_S X,
  \qquad
  (g,x)\longmapsto(gx,x),
\]
\emph{defined for every \(S'\to S\), \(g\in G(S')\), and
\(x\in X(S')\), is an epimorphism in the category of sheaves
for the topology \(\mathcal T\); equivalently, \(\Phi\) is a
covering for the topology \(\mathcal T\)
\textup{(cf.\
\hyperref[sga3:iv:proposition:4:4:3]{IV, 4.4.3})}.}

\emph{One says that \(X\) is a \emph{homogeneous space} if it
is formally homogeneous and, in addition, the morphism
\(X\to S\) is also a covering for the topology \(\mathcal T\).}

\emph{In particular, one says that \(X\) is a
\emph{formally principal homogeneous space under \(G\)} if
\(\Phi\) is an isomorphism, and that \(X\) is a
\emph{principal homogeneous bundle} \textup{(or a
\(G\)-torsor)} if \(\Phi\) is an isomorphism and, in addition,
the morphism \(X\to S\) is a covering for the topology
\(\mathcal T\) \textup{(cf.\
\hyperref[sga3:iv:definition:5:1:5]{IV, 5.1.5} and
\hyperref[sga3:iv:proposition:5:1:6:item:ii:01]{5.1.6(ii)})}.}

\medskip
\phantomsection\label{prop:sga3-VIB-10.16}%
\noindent\textbf{Proposition 10.16.---}
\emph{Retain the situation considered at the beginning of
\hyperref[par:sga3-VIB-10.0]{10.0}. Let \(j\in I\), let
\(G_j\) be an \(S_j\)-group, and let \(X_j\) be an
\(S_j\)-scheme with operator group \(G_j\). Suppose that
\(G_j\) and \(X_j\) are of finite presentation over \(S_j\).}

\emph{For}
\[
  X=X_j\times_{S_j}S
\]
\emph{to be a homogeneous space \textup{(respectively, a
principal homogeneous bundle)} under}
\[
  G=G_j\times_{S_j}S
\]
\emph{for the \textup{(fppf)} topology, it is necessary and
sufficient that there be an index \(i\geq j\) such that}
\[
  X_i=X_j\times_{S_j}S_i
\]
\emph{is a homogeneous space \textup{(respectively, a
principal homogeneous bundle)} under}
\[
  G_i=G_j\times_{S_j}S_i.
\]

Taking into account \hyperref[par:sga3-VIB-10.14]{10.14} and
EGA~IV\(_3\), 8.8.2, 8.8.3, and 8.10.5, the assertion follows
from the property of covering morphisms for the
\textup{(fppf)} topology recalled in
\hyperref[par:sga3-VIB-10.1]{10.1}.%
\footnote{\label{note:sga3-VIB-103}Editor's note: for
limit-passage properties of torsors expressed in terms of
\v{C}ech cohomology, see also SGA~4, VII, 5.7 and its
corollaries. This is developed in detail in \textup{[Ma07]}.}
